\section{Introduction}\label{sec:introduction}

Multi-turn dialogue is central to how large language models are deployed, from customer service to tutoring to long-form reasoning. Yet a persistent failure accompanies even the most fluent generation: as conversations grow longer, models pay progressively less attention to the system-prompt goals that define their task. \citet{dongre2026attention} measured significant Goal Accessibility Ratio (GAR) decline across 12 models over extended conversations; in our own evaluation, goal attention drops by over 70\% within just five turns. Crucially, this is not a representation problem: residual stream probes confirm that goal information persists (AUC\,=\,0.98). The model retains the knowledge of what it should attend to; it simply stops attending.

Existing remedies address this decay at inference time. Split-softmax \citep{li2024measuring} separates instruction from conversation tokens; StreamingLLM \citep{xiao2024streaming} retains anchor tokens; Gated Attention \citep{qiu2025gated} applies per-head sigmoid gates within a single forward pass. A natural alternative is to intervene during training: optimize goal attention directly through auxiliary losses trained alongside next-token prediction (NTP). Yet this strategy rests on an untested assumption: attention patterns learned under teacher forcing (TF) transfer to autoregressive (AR) generation. Classical exposure bias research \citep{bengio2015scheduled,lamb2016professor} has shown that TF creates a training--inference distribution mismatch for RNN sequence models, but whether this mismatch extends to auxiliary attention objectives in multi-turn transformers remains unknown.

We test this assumption and find it fails. We introduce a controlled five-way testbed (ThreadGate, +PCGrad \citep{yu2020pcgrad}, Uniform Gate, LoRA-only, and Base) designed not as competing solutions but as a diagnostic exclusion chain (\S\ref{sec:comparison-design}), where each adjacent pair isolates one candidate failure mechanism. The result is unambiguous: every variant converges to target GAR under teacher forcing, yet goal attention collapses by an order of magnitude once the model generates its own context (Figure~\ref{fig:headline}). The exclusion chain traces this collapse to exposure bias, the training--inference distribution mismatch, ruling out gradient direction conflict (PCGrad removes conflict yet \emph{worsens} AR GAR), gating mechanism (Uniform Gate matches ThreadGate), and parameter capacity (LoRA-only shows the same pattern). Inference-time $\beta$-steering confirms that the collapse reflects attention misallocation rather than representation degradation; scheduled sampling bridges the gap by exposing the model to its own generations during training.

\begin{figure*}[!ht]
  \centering
  \includegraphics[width=\textwidth]{fig1_headline.pdf}
  \caption{Exposure bias collapses auxiliary attention losses. \textbf{(a)}~Goal attention decays 85--90\% across turns under AR generation on both Qwen-2.5-7B and Llama-3.1-8B. \textbf{(b)}~Under TF evaluation all variants reach GAR$\approx$0.50, but AR generation collapses GAR by 8--17$\times$ (log scale). \textbf{(c)}~Scheduled sampling compresses the TF--AR gap from 7.8$\times$ to 1.3$\times$ (Qwen) and ${\sim}$27$\times$ to 2.9$\times$ (Llama).}
  \label{fig:headline}
\end{figure*}

Our contributions are:
\begin{enumerate}\raggedright
    \item \textbf{Exposure bias characterization.} To our knowledge, the first empirical demonstration that exposure bias collapses auxiliary attention objectives in multi-turn transformers, with an order-of-magnitude TF--AR gap reproduced across two architectures (Qwen-2.5-7B, Llama-3.1-8B) and model scales (1.5B--8B).
    \item \textbf{Diagnostic exclusion testbed.} A five-way diagnostic design, informed by seven loss-formulation iterations, that isolates the failure mechanism via systematic diagnostic exclusion, including a PCGrad \emph{modus tollens} argument ruling out gradient direction conflict ($p$\,=\,0.002, effect reverses).
    \item \textbf{Attention misallocation diagnosis.} Inference-time $\beta$-steering recovers goal attention with modest perplexity cost, evidence that the collapse is an attention routing failure, not a representational one.
    \item \textbf{Scheduled sampling as remedy.} The first application of scheduled sampling to auxiliary attention objectives in multi-turn transformers, compressing the TF--AR gap to ${\sim}$1.3$\times$ on Qwen and ${\sim}$2.9$\times$ on Llama with no perplexity degradation, robust across three seeds.
\end{enumerate}
